\chapter{平面对偶、轮廓与曲面同调}
\label{chap:phys-sm-planar-duality-contours}

一般 cycle space 只使用图的关联关系；低温轮廓和 Kramers--Wannier 对偶则需要图真正嵌入二维曲面。先在胞腔平面嵌入上建立对偶图和面边界，再把结论推广到闭可定向曲面。这样可以严格区分三件经常被混写的事实：一般偶子图只是 \(1\)-cycle；球面上的 cycle 都是面边界；正亏格曲面还保留不能由局部面边界消去的全局同调类。

\section{胞腔平面嵌入与对偶图}
\label{sec:phys-sm-planar-duality}

设连通有限多重图 \(G=(\mathsf V,\mathsf E)\) 胞腔嵌入球面 \(S^2\)，等价地可在平面上画成无交叉嵌入并把外部区域视为一个面。记面集合为 \(\mathsf F\)。胞腔条件表示每个面的内部都是开圆盘，此时 Euler 公式为
\begin{equation}
 |\mathsf V|-|\mathsf E|+|\mathsf F|=2.
 \label{eq:phys-sm-euler-plane}
\end{equation}
对偶图 \(G^*\) 的顶点对应原图面，每条原边 \(e\) 被唯一一条对偶边 \(e^*\) 横穿。原图桥对应对偶自环，原图平行边也可能产生多重对偶边，所以对偶图仍属于有限多重图范畴。

以 \(\F_2\) 为系数，引入面链空间
\begin{equation}
 C_2(G)=\F_2^{\mathsf F}
 \label{eq:phys-sm-face-chain-space}
\end{equation}
以及第二边界算子
\begin{equation}
 \partial_2:C_2(G)\longrightarrow C_1(G).
 \label{eq:phys-sm-boundary-operator-2}
\end{equation}
一个面的边界由沿该面边界出现奇数次的原边组成。因为一个面边界在每个顶点进入与离开的总次数都是偶数，
\begin{equation}
 \partial_1\partial_2=0.
 \label{eq:phys-sm-boundary-boundary}
\end{equation}
因此 \(B_1(G):=\imop\partial_2\subseteq Z_1(G)\)。

\begin{theorem}[球面胞腔嵌入中闭链都是面边界]
\label{thm:phys-sm-plane-cycle-boundary}
对连通图的球面胞腔嵌入，
\begin{equation}
 Z_1(G)=B_1(G).
 \label{eq:phys-sm-plane-z1-b1}
\end{equation}
等价地，\(\mathrm H_1(S^2;\F_2)=0\)。
\end{theorem}

\begin{derivation}[Euler 计数与面边界独立性]
分三步证明 \(Z_1(G)=B_1(G)\)。

\emph{第一步：维数相等。}由 cycle-rank 公\eqnrefs{eq:phys-sm-cycle-rank}，连通图满足
\[
 \dim Z_1=|\mathsf E|-|\mathsf V|+1.
\]
由球面胞腔嵌入的 Euler 公式 \(|\mathsf V|-|\mathsf E|+|\mathsf F|=2\)，
\[
 |\mathsf F|-1=|\mathsf E|-|\mathsf V|+1=\dim Z_1.
\]

\emph{第二步：\(\rank\partial_2\le|\mathsf F|-1\)。}每个面 \(f\) 的边界 \(\partial_2 f\) 是 \(C_1(G)\) 中的闭链。所有面边界的模二和
\[
 \sum_{f\in\mathsf F}\partial_2 f=0,
\]
因为每条原边恰好出现在两侧面的边界中，模二抵消。因此 \(\{\partial_2 f:f\in\mathsf F\}\) 线性相关，\(\rank\partial_2\le|\mathsf F|-1\)。

\emph{第三步：对偶树消元给出下界。}取原图一棵生成树 \(T\)。球面对偶中，与非树边 \(e\in\mathsf E\setminus T\) 对应的对偶边 \(e^*\) 组成对偶生成树 \(T^*\)：\(T^*\) 有 \(|\mathsf F|-1\) 条边（因为 \(|\mathsf E\setminus T|=|\mathsf E|-|\mathsf V|+1=|\mathsf F|-1\)），且 \(T^*\) 连通无环（由 \(T\) 的生成树性质对偶保证）。

删去一个参考面 \(f_0\) 后，其余 \(|\mathsf F|-1\) 个面边界 \(\{\partial_2 f:f\ne f_0\}\) 按 \(T^*\) 的叶面递归消元：取 \(T^*\) 的一个叶面 \(f_\ell\)，它通过唯一对偶树边 \(e_\ell^*\) 与其余面对偶连接。在面边界 \(\partial_2 f_\ell\) 中，\(e_\ell^*\) 对应的原边 \(e_\ell\) 恰好出现一次（因为 \(e_\ell\) 只在 \(f_\ell\) 和一个已消去的面之间），而其他待消面的边界不含 \(e_\ell\)。因此 \(\partial_2 f_\ell\) 不能被前面面边界线性表示，pivot 为 \(1\in\F_2\)。消去 \(f_\ell\) 和 \(e_\ell^*\) 后，剩余对偶图仍是更小的对偶树，归纳到单面基例。因此 \(\{\partial_2 f:f\ne f_0\}\) 线性无关，\(\rank\partial_2\ge|\mathsf F|-1\)。

结合第二步，\(\dim B_1=\rank\partial_2=|\mathsf F|-1=\dim Z_1\)。再由 \(B_1\subseteq Z_1\)（因为 \(\partial_1\partial_2=0\)），子空间维数相等即得\eqnrefs{eq:phys-sm-plane-z1-b1}。

\emph{对偶树构造的具体执行。}以四面体图 \(K_4\)（4 顶点 6 边）为例：球面嵌入给出 4 个面，cycle rank \(=6-4+1=3\)。生成树 \(T\) 取 3 条边，剩余 3 条非树边对应 3 个面（除参考面 \(f_0\)）。对偶树 \(T^*\) 是连接这 3 个面到 \(f_0\) 的对偶星图，叶面消元依次给出 3 个独立面边界，与 \(\dim Z_1=3\) 一致。每个非空偶子图（共 \(2^3-1=7\) 个）都是某面集边界的对称差。

\emph{球面假设的必要性。}若把图嵌入环面 \(T^2\)，Euler 公式改为 \(|\mathsf V|-|\mathsf E|+|\mathsf F|=0\)，于是 \(\dim Z_1=|\mathsf E|-|\mathsf V|+1=|\mathsf F|+1\)，而 \(\dim B_1\le|\mathsf F|-1\)，差值 \(\dim Z_1-\dim B_1\ge2\)。这 2 维差正是 \(\mathrm H_1(T^2;\F_2)\cong\F_2^2\) 的两个同调生成元（沿环面两个基本 cycle 的绕行），它们是闭链但不是任何面集边界。球面 \(g=0\) 时 Euler 特征 \(2\) 恰好抵消这一差值，使 \(Z_1=B_1\)。

\emph{物理后果。}\eqnrefs{eq:phys-sm-plane-z1-b1} 保证平面 Ising 的每个偶子图都能解释为某面集的边界，从而低温畴壁与高温偶子图一一对应，Kramers--Wannier 对偶在球面上成立（见\eqnrefs{eq:phys-sm-kw-finite-sphere}）。环面上额外的两个同调类对应周期/反周期边界扭转，配分函数分解为四个 spin structure 扇区，需用\eqnrefs{eq:phys-sm-torus-character-transform} 的离散 Fourier 变换处理。

\emph{对偶树消元的矩阵视角。}面边界矩阵 \(\partial_2\) 在 \(\F_2\) 上是 \(|\mathsf E|\times|\mathsf F|\) 矩阵，列对应面、行对应边。删去参考面 \(f_0\) 后的子矩阵 \(\tilde{\partial}_2\) 是 \(|\mathsf E|\times(|\mathsf F|-1)\)。对偶树消元等价于对 \(\tilde{\partial}_2\) 做高斯消元：每步选一个叶面 \(f_\ell\)，其对应列在原边 \(e_\ell\) 行有唯一 \(1\)（pivot），消去后矩阵降阶。最终得到 \(|\mathsf F|-1\) 个线性无关列，即 \(\rank\partial_2=|\mathsf F|-1\)。这一过程与生成树消元\eqnrefs{eq:phys-sm-incidence-rank} 完全对偶，体现平面图自对偶性。

\emph{Euler 公式的组合证明。}\eqnrefs{eq:phys-sm-euler-plane} 可由生成树-对偶树分解给出：边集 \(\mathsf E=T\sqcup(\mathsf E\setminus T)\)，对偶边集 \(\mathsf E^*=T^*\sqcup(\mathsf E^*\setminus T^*)\)。由 \(|T|=|\mathsf V|-1\)、\(|T^*|=|\mathsf F|-1\)、\(|\mathsf E|=|\mathsf E^*|\)，得 \(|\mathsf V|-1+|\mathsf F|-1=|\mathsf E|\)，即 Euler 公式。这一证明与同调证明互补：前者纯组合，后者通过 \(\partial_1\partial_2=0\) 与秩-零化度联系。两种证明在物理上对应高温展开（生成树）与低温畴壁（面边界）两种描述的等价性。

\emph{非平面图的失效。}对 \(K_5\)（五顶点完全图）或 \(K_{3,3}\)（三部分图），无球面胞腔嵌入（Kuratowski 定理），上述证明失效。此时 cycle space 仍由\eqnrefs{eq:phys-sm-cycle-rank} 给出，但不存在面链 \(C_2\) 使 \(Z_1=B_1\)。非平面 Ising 模型需用其他方法：高温展开仍成立（只需 cycle space），但 Kramers--Wannier 对偶失效，低温轮廓需用 cutset 或 random current formalism。这把平面性从"可选几何结构"提升为"对偶性的必要条件"。

\emph{对偶图的具体构造。}给定平面图 \(G\) 的球面嵌入，对偶图 \(G^*\) 构造如下：每个面 \(f\in\mathsf F\) 对应对偶顶点 \(f^*\)；每条原边 \(e\) 分隔两个面 \(f_1,f_2\)（可能 \(f_1=f_2\) 若 \(e\) 是桥），对应对偶边 \(e^*\) 连接 \(f_1^*,f_2^*\)。\(e\) 是桥时 \(e^*\) 是自环；\(e\) 是自环时 \(e^*\) 是桥。这一构造满足 \((G^*)^*=G\)（球面嵌入下），且 \(|\mathsf V^*|=|\mathsf F|\)、\(|\mathsf E^*|=|\mathsf E|\)、\(|\mathsf F^*|=|\mathsf V|\)，Euler 公式自洽。

\emph{方格与六角格的对偶实例。}方格 \(\mathbb Z^2\) 的对偶仍是方格（自对偶），这是 Kramers--Wannier 自对偶点 \(\sinh(2K_c)=1\) 存在的几何基础。六角格的对偶是三角格，对偶耦合 \(K_{\rm hex}^*\) 与 \(K_{\rm tri}^*\) 满足 \(\sinh(2K_{\rm hex})\sinh(2K_{\rm tri})=2/\sqrt{3}\)（非对称常数来自不同配位数）。三角格自对偶（对偶仍是三角格），自对偶点 \(\sinh(2K_c)=2/\sqrt{3}\)。这些具体对偶关系在相变临界点计算中直接应用，给出不同晶格 Ising 模型的精确 \(T_c\)。

\emph{面边界矩阵的秩-零空间结构。}\(\partial_2:C_2\to C_1\) 的像 \(B_1=\imop\partial_2\) 是面边界空间，核 \(\ker\partial_2\) 是"面链中边界为零"的组合。球面上 \(\ker\partial_2\) 一维，由所有面的和 \(\sum_{f\in\mathsf F}f\) 生成（其边界为 \(\sum\partial_2 f=0\)）。环面上 \(\ker\partial_2\) 仍一维（同一全局关系），但 \(\dim B_1=|\mathsf F|-1\) 不变，\(\dim Z_1\) 增加 2，差值进入 \(\mathrm H_1\)。这一核-像结构是链复形正合性的体现，连接同调维数与 Euler 特征。

\emph{物理建模的完整链条。}上述代数结构在 Ising 模型中按以下链条应用：
\begin{equation*}
 \text{图 }G
 \xrightarrow{\;\text{胞腔嵌入}\;}
 \text{面链 }C_2
 \xrightarrow{\;\partial_2\;}
 B_1\subseteq Z_1
 \xrightarrow{\;Z_1/B_1\;}
 \mathrm H_1
 \xrightarrow{\;\text{spin structure}\;}
 \text{配分函数分解}.
\end{equation*}
球面 \(\mathrm H_1=0\) 给出单一配分函数（Kramers--Wannier 对偶直接成立）；环面 \(\mathrm H_1\cong\F_2^2\) 给出四扇区分解（需离散 Fourier 变换）；高亏格 \(g\) 给出 \(2^{2g}\) 扇区。每一步的代数结构（cycle rank、Euler 特征、同调维数）都对应物理量（偶子图数、面边界独立数、拓扑扇区数），把抽象代数与具体统计力学严格联系起来。
\end{derivation}

\section{固定边界下的低温轮廓}
\label{sec:phys-sm-contours}

考虑单连通平面有限区域中的最近邻 Ising 模型，并把外边界自旋固定为 \(+1\)。对任意自旋组态 \(\sigma\)，若原边 \(e=\{x,y\}\) 满足 \(\sigma_x\neq\sigma_y\)，就在对偶边 \(e^*\) 上放置畴壁。所有畴壁组成 \(\Gamma(\sigma)\subseteq\mathsf E^*\)。沿任意原图面走一周，自旋翻转次数必须为偶数，因此
\begin{equation}
 \partial_1^*\Gamma(\sigma)=0.
 \label{eq:phys-sm-domain-wall-even}
\end{equation}
低温畴壁因此是对偶图上的闭 \(1\)-链。

\begin{theorem}[固定加边界下的轮廓--组态一一对应]
\label{thm:phys-sm-spin-contour-bijection}
在单连通平面有限区域上固定外边界为 \(+\) 后，自旋组态与满足局部偶度条件且不违反固定边界的对偶闭链一一对应。
\end{theorem}

\begin{derivation}[由闭轮廓唯一重建自旋]
分三步证明自旋组态与对偶闭链的一一对应。

\emph{第一步：自旋组态产生闭链。}给定自旋组态时，\eqnrefs{eq:phys-sm-domain-wall-even} 已说明畴壁 \(\Gamma(\sigma)\) 满足 \(\partial_1^*\Gamma=0\)，必为闭链。

\emph{第二步：闭链重建自旋（良定性）。}给定允许的对偶闭链 \(\Gamma\)，取任意顶点 \(x\)，选择从外边界参考点 \(x_\partial\) 到 \(x\) 的原图路径 \(P_x\)，定义
\[
 \sigma_x=(-1)^{I(P_x,\Gamma)},
\]
其中 \(I(P_x,\Gamma)\in\F_2\) 是路径与对偶轮廓的模二交数，外边界参考自旋取 \(+1\)。

需证 \(\sigma_x\) 与路径选择无关。若 \(P_x\) 与 \(P_x'\) 是两条从 \(x_\partial\) 到 \(x\) 的路径，则它们的模二和
\[
 C=P_x+P_x'
\]
是原图闭链（起点和终点在模二中相消）。在单连通平面区域中，\(C\) 是若干面边界的模二和：\(C=\sum_{f\in S}\partial_2 f\)（某面集 \(S\)）。而闭对偶链 \(\Gamma\) 穿过任意面边界 \(\partial_2 f\) 的模二交数为零，因为 \(\Gamma\) 进入面 \(f\) 的次数等于离开次数。因此
\[
 I(C,\Gamma)=\sum_{f\in S}I(\partial_2 f,\Gamma)=0\pmod2,
\]
故
\[
 I(P_x,\Gamma)-I(P_x',\Gamma)=I(P_x+P_x',\Gamma)=I(C,\Gamma)=0\pmod2.
\]
\(\sigma_x\) 良定。

\emph{第三步：重建组态产生原来的 \(\Gamma\) 且一一。}对原边 \(e=\{x,y\}\)，取路径 \(P_x\) 到 \(x\)，则 \(P_y=P_x+e\) 到 \(y\)。若对偶边 \(e^*\in\Gamma\)，则
\[
 \sigma_y=(-1)^{I(P_y,\Gamma)}=(-1)^{I(P_x,\Gamma)+1}=-\sigma_x,
 \qquad \sigma_x\sigma_y=-1.
\]
若 \(e^*\notin\Gamma\)，则 \(\sigma_y=\sigma_x\)，\(\sigma_x\sigma_y=+1\)。因此重建组态恰好产生原来的 \(\Gamma\)。固定外边界为 \(+\) 去掉了整体自旋翻转 \(\sigma\to-\sigma\) 的二重性（翻转使所有 \(\sigma_x\sigma_y\) 不变），因此对应是一一的。

\emph{单连通性的具体作用。}良定性证明中关键步骤 \(C=\sum_{f\in S}\partial_2 f\) 依赖原图闭链 \(C\) 在球面上都是面边界，即\eqnrefs{eq:phys-sm-plane-z1-b1}。环面上存在非边界的闭链（如沿基本 cycle \(a_x\) 的绕行），若 \(P_x+P_x'\) 恰好是这种同调非平凡闭链，则 \(I(C,\Gamma)\) 可能等于 \(1\)，\(\sigma_x\) 不良定。此时需额外指定同调扇区（绕行数 \((w_x,w_y)\in\F_2^2\)）才能恢复良定自旋，这正是环面四个 spin structure 的来源。

\emph{畴壁能量的具体计算。}对一组闭畴壁 \(\Gamma\)，全加组态 \(\sigma\equiv+1\) 的能量为 \(E_0=-\sum_e J_e\)。每条畴壁边 \(e\) 使 \(\sigma_x\sigma_y=-1\)，局部能量从 \(-J_e\) 变为 \(+J_e\)，增量 \(2J_e\)。总能量
\begin{equation*}
 E(\Gamma)=E_0+2\sum_{e^*\in\Gamma}J_e
 =-\sum_e J_e+2\sum_{e^*\in\Gamma}J_e,
\end{equation*}
故 Boltzmann 权重 \(\ee^{-\beta E(\Gamma)}=\ee^{\sum_e K_e}\prod_{e^*\in\Gamma}\ee^{-2K_e}\)，给出\eqnrefs{eq:phys-sm-low-temp-z}。固定边界消除整体翻转二重性，故求和遍历所有允许闭链而不需额外因子 \(2\)。

\emph{Peierls 估计的轮廓计数。}在二维方格中，围住原点、长度 \(\ell\) 的对偶闭链数 \(N_\ell\le\ell\cdot 3^{\ell-1}\)（每步至多 3 个非回溯选择，长度 \(\ell\) 的起点有 \(\ell\) 个可能位置）。当 \(\mu\ee^{-2\beta J_*}<1\) 即 \(T<T_{\rm Peierls}=2J_*/\log(3\mu)\) 时，\eqnrefs{eq:phys-sm-peierls-bound-general} 收敛，\(\mathbb P_+(\sigma_0=-1)<1/2\)，自发磁化非零。这给出二维 Ising 低温相变的定性证明，临界温度 \(T_c\) 的精确值需由 Kramers--Wannier 自对偶\eqnrefs{eq:phys-sm-square-critical-k} 或 Onsager 谱确定。

\emph{低温展开与高温展开的比较。}对二维 \(L\times L\) 方格铁磁 Ising，两种展开在不同温度区间收敛：
\begin{itemize}
 \item 低温展开（畴壁）：\(Z=2\ee^{L^2 K}\sum_\Gamma\ee^{-2K|\Gamma|}\)，收敛条件 \(\ee^{-2K}\ll1\) 即 \(K\gg1\)（\(T\ll T_c\)）。首项为全加组态 \(2\ee^{L^2 K}\)，次项为最小畴壁（单 plaquette）\(4L^2\ee^{L^2 K}\ee^{-8K}\)。
 \item 高温展开（偶子图）：\(Z=2^{L^2}(\cosh K)^{2L^2}\sum_{\mathscr A\in Z_1}(\tanh K)^{|\mathscr A|}\)，收敛条件 \(\tanh K\ll1\) 即 \(K\ll1\)（\(T\gg T_c\)）。首项为空偶子图 \(2^{L^2}(\cosh K)^{2L^2}\)，次项为单 plaquette \(L^2(\tanh K)^4\)。
\end{itemize}
两者通过 Kramers--Wannier 对偶 \(K\leftrightarrow K^*\)（\(\sinh 2K\sinh 2K^*=1\)）严格联系，临界点 \(K_c\) 恰为两个展开收敛半径的边界。低温展开直接显示自发磁化（全加组态主导），高温展开直接显示解析性（\(\tanh K\) 小时 \(\mathcal Z_G\) 解析），两者在 \(K_c\) 处汇合给出完整相图。

\emph{固定边界条件的物理意义。}固定外边界为 \(+\) 选取纯铁磁相，对应热力学极限 \(h\to0^+\)（正场极限）。若固定为 \(-\)，得到纯铁磁相的翻转 \(h\to0^-\)。自由边界条件对应 \(h=0\) 对称混合相，低温下两相共存。周期边界条件（环面）需额外指定同调扇区，四个扇区在热力学极限下都趋于同一纯相（平凡扇区主导），但有限尺寸修正依赖扇区选择。这一边界条件敏感性是相变非局域性的体现：低温下局部边界选择能决定整个体积的相。
\end{derivation}

这个定理的单连通性不能删除。周期边界上闭链可能具有非零绕行数，无法仅靠“从外部区域开始翻转”全局定义自旋；一般非平面图甚至没有对偶轮廓包围区域的概念。

\section{闭可定向曲面上的第一同调}
\label{sec:phys-sm-homology}

设图给出闭可定向 genus \(g\) 曲面 \(\Sigma_g\) 的胞腔分解。仍有链复形
\[
 C_2\xrightarrow{\partial_2}C_1\xrightarrow{\partial_1}C_0,
 \qquad \partial_1\partial_2=0,
\]
第一同调定义为
\begin{equation}
 \mathrm H_1(\Sigma_g;\F_2)
 =\kerop\partial_1/\imop\partial_2.
 \label{eq:phys-sm-first-homology}
\end{equation}
它把闭链按“只相差局部面边界”分类。

对闭可定向曲面，Euler 特征为 \(\chi(\Sigma_g)=2-2g\)。连通性给出 \(\dim\ker\partial_1=|\mathsf E|-|\mathsf V|+1\)。另一方面，全部面边界只有一条总和为零的关系，因此 \(\rank\partial_2=|\mathsf F|-1\)。于是
\begin{equation}
 \boxed{
 \dim_{\F_2}\mathrm H_1(\Sigma_g;\F_2)=2g.
 }
 \label{eq:phys-sm-genus-h1-dimension}
\end{equation}
因此共有 \(2^{2g}\) 个模二同调扇区。球面 \(g=0\) 恢复 \(\mathrm H_1=0\)，环面 \(g=1\) 则给出
\begin{equation}
 \mathrm H_1(T^2;\F_2)\cong\F_2^2.
 \label{eq:phys-sm-torus-h1}
\end{equation}
若考虑圆柱 \(S^1\times[0,1]\)，有一个周期方向，因此
\begin{equation}
 \mathrm H_1(S^1\times[0,1];\F_2)\cong\F_2.
 \label{eq:phys-sm-cylinder-h1}
\end{equation}

\begin{derivation}[由 Euler 特征得到 \eqnrefs{eq:phys-sm-torus-h1} 中的 \(2g\)]
闭曲面 \(\Sigma_g\) 的亏格 \(g\) 通过 Gauss--Bonnet 编码几何、通过 Euler 特征编码组合。在任意胞腔剖分下，Euler 特征
\begin{equation*}
 \chi(\Sigma_g)=|\mathsf V|-|\mathsf E|+|\mathsf F|=2-2g
\end{equation*}
是拓扑不变量，与剖分选取无关。另一方面，链复形
\begin{equation*}
 0\xrightarrow{\partial_2}C_1(G)\xrightarrow{\partial_1}C_0(G)\to 0
\end{equation*}
给出两个维数恒等式。第一，闭链空间 \(Z_1=\ker\partial_1\) 的维数由 cycle-rank 公\eqnrefs{eq:phys-sm-cycle-rank}给出（连通图 \(c(G)=1\)）：
\begin{equation*}
 \dim Z_1=|\mathsf E|-|\mathsf V|+1.
\end{equation*}
第二，边缘空间 \(B_1=\operatorname{im}\partial_2\) 的维数等于面边界矩阵的秩。每个面的边界是闭 \(1\)-链，但所有面边界的模二和为零（每条内部边恰出现在两个面的边界中），故至少有一个线性关系。球面情形这唯一的关系穷尽了所有依赖，但高亏格面还需要额外的同调生成元。对任意亏格，
\begin{equation*}
 \dim B_1=|\mathsf F|-1,
\end{equation*}
因为面边界张成的空间维数恰好是面数减去这个全局关系。

将两式代入同调定义 \(\mathrm H_1=Z_1/B_1\)，\(\dim\mathrm H_1=\dim Z_1-\dim B_1\)：
\begin{align*}
 \dim\mathrm H_1
 &=\bigl(|\mathsf E|-|\mathsf V|+1\bigr)-\bigl(|\mathsf F|-1\bigr)\\
 &=|\mathsf E|-|\mathsf V|-|\mathsf F|+2\\
 &=-\bigl(|\mathsf V|-|\mathsf E|+|\mathsf F|\bigr)+2\\
 &=-(2-2g)+2=2g.
\end{align*}
因此 \(\mathrm H_1(\Sigma_g;\F_2)\cong\F_2^{2g}\)，即\eqnrefs{eq:phys-sm-torus-h1}。球面 \(g=0\) 给出 \(\dim\mathrm H_1=0\)，环面 \(g=1\) 给出 \(\dim\mathrm H_1=2\)（四个同调扇区），双环面 \(g=2\) 给出 \(\dim\mathrm H_1=4\)（十六个扇区）。

\emph{亏格的几何与组合双重刻画。}Gauss--Bonnet \(\int_{\Sigma_g}K\,\dd A=2\pi\chi=4\pi(1-g)\) 把亏格与曲率积分联系；胞腔剖分 \(|\mathsf V|-|\mathsf E|+|\mathsf F|=2-2g\) 把同一亏格与离散计数联系。两者等价性来自剖分光滑化与曲率集中在顶点的极限过程。物理上，Ising 模型在连续曲面上的临界行为只依赖 \(g\)，与具体剖分选取无关——这是拓扑不变性在统计力学中的体现。

\emph{环面 \(g=1\) 的四个扇区。}取 \(L\times L\) 方格周期边界（环面），\(\dim\mathrm H_1=2\)，四个同调类由绕行数 \((w_x,w_y)\in\F_2^2\) 标记：\((0,0)\) 平凡同调（可缩闭链），\((1,0)\) 沿 \(x\)-cycle 绕行一次，\((0,1)\) 沿 \(y\)-cycle，\((1,1)\) 同时沿两个方向。每个扇区对配分函数贡献 \(Z_{w_x,w_y}\)，总配分函数 \(Z=\sum_{w_x,w_y}Z_{w_x,w_y}\)。Kac--Ward 行列式或 transfer matrix 分别计算四个扇区，再由\eqnrefs{eq:phys-sm-torus-character-inverse} 反变换组合。Onsager 自由能取平凡扇区 \((0,0)\) 主导项，其余扇区在热力学极限下指数 suppressed。

\emph{高亏格曲面的物理应用。}双环面 \(g=2\)（如八角格嵌入）给出 \(2^4=16\) 个同调扇区，配分函数分解为 16 项。一般亏格 \(g\) 曲面上 Ising 临界行为仍由局部耦合决定（普适类不变），但有限尺寸修正依赖 \(g\)：自由能 \(f_L=f_\infty+\alpha(g)/L^2+\cdots\)，\(\alpha(g)\) 通过同调维数 \(2g\) 进入。这把拓扑不变量与有限尺寸标度理论直接联系起来，是 Ising 普适类在曲面族上稳定性的体现。

\emph{同调维数 \(2g\) 的几何来源。}\(2g\) 个生成元对应曲面上 \(g\) 个"环柄"各自贡献的两个基本 cycle：经线（绕环柄短轴）与纬线（绕环柄长轴）。环面 \(g=1\) 有一对 \((a_x,a_y)\)；双环面 \(g=2\) 有两对 \((a_1,b_1),(a_2,b_2)\)，四条独立 cycle。亏格 \(g\) 曲面的基本群 \(\pi_1(\Sigma_g)\) 有 \(2g\) 个生成元满足单一关系 \(\prod_{i=1}^g[a_i,b_i]=1\)，其 Abel 化给出 \(\mathrm H_1\cong\mathbb Z^{2g}\)，模二化得 \(\F_2^{2g}\)。

\emph{Ising 配分函数的扇区分解。}亏格 \(g\) 曲面上 Ising 配分函数分解为 \(2^{2g}\) 个同调扇区贡献：
\begin{equation*}
 Z_{\Sigma_g}=\sum_{\mathbf w\in\F_2^{2g}}Z_{\mathbf w},
 \qquad \mathbf w=(w_1,\ldots,w_{2g}),
\end{equation*}
每个 \(Z_{\mathbf w}\) 固定绕行数 \(\mathbf w\)。对应的扭转配分函数 \(Z^{\boldsymbol\upsilon}\)（\(\boldsymbol\upsilon\in\F_2^{2g}\)）通过 \(2^{2g}\times2^{2g}\) 离散 Fourier 变换联系：
\begin{equation*}
 Z^{\boldsymbol\upsilon}=\sum_{\mathbf w}(-1)^{\boldsymbol\upsilon\cdot\mathbf w}Z_{\mathbf w},
 \qquad
 Z_{\mathbf w}=2^{-2g}\sum_{\boldsymbol\upsilon}(-1)^{\boldsymbol\upsilon\cdot\mathbf w}Z^{\boldsymbol\upsilon}.
\end{equation*}
Kac--Ward/Pfaffian 方法直接计算 \(Z^{\boldsymbol\upsilon}\)（每个扭转对应一个 spin structure），再反变换得到 \(Z_{\Sigma_g}\)。热力学极限下平凡扇区 \(\mathbf w=0\) 主导，其余扇区指数 suppressed，但有限尺寸下所有扇区贡献相当。

\emph{普适类与拓扑独立性。}Ising 临界指数 \((\alpha,\beta,\gamma,\delta,\nu,\eta)=(0,1/8,7/4,15,1,1/4)\) 在所有亏格 \(g\) 曲面上相同，因为临界行为由局部涨落（关联长度 \(\xi\gg1\)）决定，全局拓扑只影响有限尺寸修正。这体现普适类原则：临界指数只依赖空间维数 \(d\) 与对称性（\(\mathbb Z_2\)），不依赖晶格细节（方格、六角、三角）或全局拓扑（球面、环面、高亏格）。拓扑信息通过 \(\dim\mathrm H_1=2g\) 进入有限尺寸标度函数，但不改变临界指数本身。
\end{derivation}

\begin{derivation}[Poincaré 对偶与曲面上的辛相交形式]
闭可定向曲面 \(\Sigma_g\) 上，第一同调 \(\mathrm H_1(\Sigma_g;\F_2)\cong\F_2^{2g}\) 与第一上同调 \(\mathrm H^1(\Sigma_g;\F_2)\cong\F_2^{2g}\) 自然同构。Poincaré 对偶给出更精确的结论：相交配对给出非退化双线性形式，使 \(\mathrm H_1\) 成为辛向量空间。分四步建立。

\emph{第一步：相交配对的定义与良定性。}对两个闭 \(1\)-链 \(\mathscr A,\mathscr B\in Z_1(G)\)，取横截代表元后定义模二交数
\begin{equation*}
  I(\mathscr A,\mathscr B)=|\mathscr A\cap\mathscr B|\pmod2.
\end{equation*}
先证 \(I\) 只依赖同调类。若 \(\mathscr A'=\mathscr A+\partial_2 f\)（相差一个面边界），则 \(\mathscr A'\) 与 \(\mathscr B\) 比 \(\mathscr A\) 多了 \(\partial_2 f\) 与 \(\mathscr B\) 的交。闭链 \(\mathscr B\) 进入面 \(f\) 的次数等于离开次数，故 \(I(\partial_2 f,\mathscr B)=0\pmod2\)，即 \(I(\mathscr A',\mathscr B)=I(\mathscr A,\mathscr B)\)。同理对第二个变量。因此 \(I\) 诱导
\begin{equation*}
  \omega:\mathrm H_1(\Sigma_g;\F_2)\times\mathrm H_1(\Sigma_g;\F_2)\longrightarrow\F_2.
\end{equation*}

\emph{第二步：交错性（特征 \(2\) 下的对称性）。}在 \(\F_2\) 上，\(\omega(\alpha,\beta)=|\alpha\cap\beta|\pmod2=|\beta\cap\alpha|\pmod2=\omega(\beta,\alpha)\)，故 \(\omega\) 对称。但曲面的相交形式在整数系数下是交错的：\(I_\mathbb Z(\alpha,\beta)=-I_\mathbb Z(\beta,\alpha)\)。模二化后交错性变为对称性，\(\omega(\alpha,\alpha)=0\)（自交数为偶：闭链可微扰使自交成对出现）。因此 \(\omega\) 是 \(\F_2\)-辛形式：对称且自交为零。

\emph{第三步：标准辛基的构造。}亏格 \(g\) 曲面有 \(g\) 个环柄，每环柄给出一对基本 cycle \((a_i,b_i)\)，\(i=1,\ldots,g\)。它们的相交数为
\begin{equation*}
  \omega(a_i,a_j)=0,\quad\omega(b_i,b_j)=0,\quad\omega(a_i,b_j)=\delta_{ij}.
\end{equation*}
前两式因为同类 cycle 可取不相交代表；第三式因为 \(a_i\) 与 \(b_i\) 在第 \(i\) 环柄上横截相交一次，不同环柄的 cycle 不相交。在基 \(\{a_1,b_1,\ldots,a_g,b_g\}\) 下，\(\omega\) 的矩阵为分块直和
\begin{equation*}
  \Omega=\bigoplus_{i=1}^g\begin{pmatrix}0&1\\1&0\end{pmatrix},
\end{equation*}
即 \(2g\times2g\) 非退化矩阵（\(\det\Omega=1\)）。因此 \(\omega\) 非退化。

\emph{第四步：Poincaré 对偶与完美配对。}非退化双线性形式 \(\omega\) 诱导线性映射
\begin{equation*}
  \mathrm{PD}:\mathrm H_1(\Sigma_g;\F_2)\longrightarrow\mathrm H^1(\Sigma_g;\F_2),\quad\mathrm{PD}(\alpha)=\omega(\alpha,\cdot).
\end{equation*}
\(\omega\) 非退化等价于 \(\mathrm{PD}\) 为同构。这给出曲面上的 Poincaré 对偶：\(\mathrm H_1\cong\mathrm H^1\)，且对偶由相交形式自然实现。结合\eqnrefs{eq:phys-sm-genus-h1-dimension}，\(\dim\mathrm H^1=2g\)，与 \(\dim\mathrm H_1=2g\) 一致。

\emph{辛结构对配分函数分解的约束。}\(\omega\) 的辛结构直接约束 Ising 配分函数的扇区分解。扭转 character \(\boldsymbol\upsilon\in\mathrm H^1\cong\mathrm H_1\)（通过 Poincaré 对偶）对应一个同调类 \(\alpha_\upsilon\)。两个扭转 \(\boldsymbol\upsilon,\boldsymbol\upsilon'\) 的"叠加"对应辛和 \(\alpha_\upsilon+\alpha_{\upsilon'}\)。Kac--Ward Pfaffian 在扭转 \(\boldsymbol\upsilon\) 下的符号由 \(\omega(\alpha_\upsilon,\alpha_{\upsilon'})\) 控制——这是 spin structure 的 Arf 不变量理论的组合基础。具体地，\(2^{2g}\) 个扭转中，Arf 不变量为 \(0\) 的有 \(2^{g-1}(2^g+1)\) 个（偶 spin structure），为 \(1\) 的有 \(2^{g-1}(2^g-1)\) 个（奇 spin structure）。环面 \(g=1\) 给出 \(3\) 偶 \(1\) 奇，与 Kac--Ward 四扇区中三个 Pfaffian 非零、一个为零一致。

\emph{整数系数与模二化的关系。}整数同调 \(\mathrm H_1(\Sigma_g;\mathbb Z)\cong\mathbb Z^{2g}\) 上的相交形式是标准辛矩阵 \(\bigoplus_i\bigl(\begin{smallmatrix}0&1\\-1&0\end{smallmatrix}\bigr)\)。模二化把 \(-1\) 映为 \(1\)，交错变为对称。万有系数定理给出 \(\mathrm H_1(\Sigma_g;\F_2)\cong\mathrm H_1(\Sigma_g;\mathbb Z)\otimes\F_2\)，因此模二同调的辛结构是整数辛结构的约化。物理上，Ising 模型的 \(\mathbb Z_2\) 对称性自然选择模二系数，但更精细的 spin structure 分类需要整数系数的辛结构（区分偶/奇 spin structure）。辛结构的维数公式为
\begin{equation*}
\dim_{\F_2}\mathrm H_1(\Sigma_g;\F_2)=2g,\qquad |\mathrm{Sp}(2g,\F_2)|=2^{g^2}\prod_{k=1}^g(2^{2k}-1).
\end{equation*}






\end{derivation}

\section{环面交数、扭转与离散 Fourier 变换}

在环面上选取两个基本 \(1\)-cycle \(a_x,a_y\) 及其横截割线 \(\ell_x,\ell_y\)。对闭链 \(\mathscr A\) 定义模二交数
\begin{equation}
 w_x(\mathscr A)=I(\mathscr A,\ell_x),
 \qquad
 w_y(\mathscr A)=I(\mathscr A,\ell_y),
 \label{eq:phys-sm-winding-parities}
\end{equation}
其中 \(I\in\F_2\)。面边界与闭割线的交数为零，因此 \((w_x,w_y)\) 只依赖同调类。

\begin{derivation}[交数为何给出完整的环面同调坐标]
环面 \(T^2\) 的第一同调已由\eqnrefs{eq:phys-sm-torus-h1}给出 \(\mathrm H_1(T^2;\F_2)\cong\F_2^2\)，因此任一同调类 \([\mathscr A]\) 可唯一写成基本 cycle \(a_x,a_y\) 的模二线性组合
\[
 [\mathscr A]=c_x[a_x]+c_y[a_y],
 \qquad c_x,c_y\in\F_2.
\]
取与 \(a_x,a_y\) 横截对偶的割线 \(\ell_x,\ell_y\)，使基本交数为
\[
 I(a_x,\ell_x)=1,
 \quad I(a_x,\ell_y)=0,
 \quad
 I(a_y,\ell_x)=0,
 \quad I(a_y,\ell_y)=1.
\]
由面边界与闭割线的交数为零，交数 \(I(\mathscr A,\ell)\) 只依赖同调类 \([\mathscr A]\) 而不依赖代表元 \(\mathscr A\)。于是定义映射
\[
 \Phi:\mathrm H_1(T^2;\F_2)\longrightarrow\F_2^2,
 \qquad
 \Phi([\mathscr A])=
 \bigl(I(\mathscr A,\ell_x),I(\mathscr A,\ell_y)\bigr).
\]
先证 \(\Phi\) 是线性映射。若 \([\mathscr A],[\mathscr B]\in\mathrm H_1(T^2;\F_2)\)，取代表元使 \(\mathscr A+\mathscr B\) 为模二链和，则
\[
 \begin{aligned}
 \Phi([\mathscr A]+[\mathscr B])
 &=\bigl(I(\mathscr A+\mathscr B,\ell_x),I(\mathscr A+\mathscr B,\ell_y)\bigr)\\
 &=\bigl(I(\mathscr A,\ell_x)+I(\mathscr B,\ell_x),I(\mathscr A,\ell_y)+I(\mathscr B,\ell_y)\bigr)\\
 &=\Phi([\mathscr A])+\Phi([\mathscr B]),
 \end{aligned}
\]
其中用到模二交数在第一个变量上的双线性。
再算 \(\Phi\) 在基本同调类上的值。由 \([a_x]=1\cdot[a_x]+0\cdot[a_y]\)，
\[
 \Phi([a_x])
 =\bigl(I(a_x,\ell_x),I(a_x,\ell_y)\bigr)
 =(1,0),
\]
由 \([a_y]=0\cdot[a_x]+1\cdot[a_y]\)，
\[
 \Phi([a_y])
 =\bigl(I(a_y,\ell_x),I(a_y,\ell_y)\bigr)
 =(0,1).
\]
\((1,0),(0,1)\) 是 \(\F_2^2\) 的标准基，因此 \(\Phi\) 把 \(\mathrm H_1(T^2;\F_2)\) 的生成元映到 \(\F_2^2\) 的生成元，必为满射。又由\eqnrefs{eq:phys-sm-torus-h1}两端都是二维 \(\F_2\)-向量空间，满射在同维有限维向量空间之间必为同构，故 \(\Phi\) 是同构。
最后把 \(\Phi\) 作用到一般同调类上。由 \([\mathscr A]=c_x[a_x]+c_y[a_y]\) 与 \(\Phi\) 的线性，
\[
 \Phi([\mathscr A])
 =c_x\Phi([a_x])+c_y\Phi([a_y])
 =c_x(1,0)+c_y(0,1)
 =(c_x,c_y).
\]
因此两个交数 \(\bigl(I(\mathscr A,\ell_x),I(\mathscr A,\ell_y)\bigr)\) 恰好给出 \([\mathscr A]\) 在基 \(\{[a_x],[a_y]\}\) 下的坐标，即\eqnrefs{eq:phys-sm-winding-parities}中的 \((w_x,w_y)\) 是 \(\mathrm H_1(T^2;\F_2)\) 的完整同调坐标，而非附加的方便标签。

\emph{物理意义。}交数作为同调坐标的物理意义是"环绕数完全刻画环面上的拓扑类"：在二维环面 \(T^2\) 上，任何闭曲线的同调类由它与两个基本循环的交数唯一确定，这两个交数正是物理学中的环绕数 \((w_x,w_y)\)。这把"粒子在环面上绕了多少圈"这一拓扑信息编码为 \(\mathbb F_2\) 上的线性代数——交数模 2 给出同调群的坐标。在 Kitaev 自旋模型中，任意子（anyon）的熔接统计由环绕数决定：m-任意子绕 e-任意子一周获得 \(-1\) 相位，正是交数的物理体现。这一同调坐标的完备性保证：两个闭曲线同伦当且仅当它们的交数相同，使拓扑量子计算中的编织操作可完全由交数矩阵描述。
\end{derivation}

沿某个周期方向插入反周期 Ising 割线，等价于把所有穿过该割线的耦合变号。高温闭链每穿过一次割线就获得 \(-1\)，所以边界扭转定义了 \(\mathrm H_1(T^2;\F_2)\) 上的 character，也就是 \(\mathrm H^1(T^2;\F_2)\) 元素。若 \(Z_{w_x,w_y}\) 是固定同调扇区的图和，\(Z^{\upsilon_x,\upsilon_y}\) 是固定扭转 \(\upsilon_x,\upsilon_y\in\F_2\) 的配分函数，则二者满足
\begin{equation}
 Z^{\upsilon_x,\upsilon_y}
 =\sum_{w_x,w_y\in\F_2}
 (-1)^{\upsilon_xw_x+\upsilon_yw_y}
 Z_{w_x,w_y},
 \label{eq:phys-sm-torus-character-transform}
\end{equation}
以及逆变换
\begin{equation}
 Z_{w_x,w_y}
 =\frac14\sum_{\upsilon_x,\upsilon_y\in\F_2}
 (-1)^{\upsilon_xw_x+\upsilon_yw_y}
 Z^{\upsilon_x,\upsilon_y}.
 \label{eq:phys-sm-torus-character-inverse}
\end{equation}
因此“同调扇区”和“边界 spin structure”不是逐项相同的对象：前者属于 \(\mathrm H_1\)，后者在选择参考结构后由 \(\mathrm H^1\) 的 character 标记，二者通过有限 Fourier 变换相连。后续 Kac--Ward/Pfaffian 与 Jordan--Wigner 两种语言都会再次出现这组四重结构。

\begin{derivation}[Dehn 扭转与映射类群的辛表示]
曲面的映射类群 \(\mathrm{MCG}(\Sigma_g)=\pi_0\mathrm{Diff}^+(\Sigma_g)\) 作用在 \(\mathrm H_1(\Sigma_g;\mathbb Z)\cong\mathbb Z^{2g}\) 上，保持相交形式不变，给出辛表示 \(\rho:\mathrm{MCG}(\Sigma_g)\to\mathrm{Sp}(2g,\mathbb Z)\)。这一表示是理解 Ising 模型在曲面变换下行为的关键。分四步建立。

\emph{第一步：Dehn 扭转的定义与同调作用。}沿简单闭曲线 \(\gamma\) 的 Dehn 扭转 \(T_\gamma\) 是曲面微分同胚：在 \(\gamma\) 的管状邻域 \(S^1\times[0,1]\) 上，\(T_\gamma(\theta,t)=(\theta+2\pi t,t)\)，邻域外为恒等。对任意闭链 \(\alpha\)，\(T_\gamma\) 的同调作用为
\begin{equation*}
  T_{\gamma,*}(\alpha)=\alpha+I_\mathbb Z(\alpha,\gamma)\,\gamma,
\end{equation*}
其中 \(I_\mathbb Z\) 是整数相交数。证明：在 \(\gamma\) 邻域外 \(\alpha\) 不变；在邻域内，\(\alpha\) 每穿过 \(\gamma\) 一次，被扭转 \(T_\gamma\) 沿 \(\gamma\) 方向推一段，净效果是增加 \(I_\mathbb Z(\alpha,\gamma)\) 个 \(\gamma\) 的拷贝。

\emph{第二步：辛性质的验证。}对两个闭链 \(\alpha,\beta\)，
\begin{equation*}
  \begin{aligned}
  I_\mathbb Z(T_{\gamma,*}\alpha,T_{\gamma,*}\beta)
  &=I_\mathbb Z(\alpha+I(\alpha,\gamma)\gamma,\;\beta+I(\beta,\gamma)\gamma)\\
  &=I(\alpha,\beta)+I(\alpha,\gamma)I(\gamma,\beta)+I(\gamma,\alpha)I(\beta,\gamma)+0\\
  &=I(\alpha,\beta)+I(\alpha,\gamma)I(\gamma,\beta)-I(\alpha,\gamma)I(\gamma,\beta)\\
  &=I(\alpha,\beta).
  \end{aligned}
\end{equation*}
第三行用了 \(I(\gamma,\alpha)=-I(\alpha,\gamma)\)，第四行两项抵消。因此 \(T_{\gamma,*}\) 保持相交形式，\(\rho(T_\gamma)\in\mathrm{Sp}(2g,\mathbb Z)\)。

\emph{第三步：标准生成元与辛矩阵。}取标准辛基 \(\{a_1,b_1,\ldots,a_g,b_g\}\)。沿 \(a_i\) 的 Dehn 扭转 \(T_{a_i}\) 的作用为
\begin{equation*}
  T_{a_i,*}(a_j)=a_j,\qquad T_{a_i,*}(b_j)=b_j+\delta_{ij}a_j,
\end{equation*}
对应辛矩阵为单位矩阵加上 \(a_i\) 行 \(b_i\) 列的 \(1\)。沿 \(b_i\) 类似。这 \(2g\) 个 Dehn 扭转生成（事实上连同少量额外扭转）整个映射类群（Dehn--Lickorish 定理），其辛表示生成 \(\mathrm{Sp}(2g,\mathbb Z)\)。

\emph{第四步：表示的核与 Torelli 群。}\(\rho\) 的核 \(\ker\rho=\mathrm{I}(\Sigma_g)\) 是 Torelli 群——同调作用下为恒等的微分同胚。\(g\ge2\) 时 Torelli 群非平凡（沿分离曲线的 Dehn 扭转在 \(\mathrm H_1\) 上作用平凡），\(\rho\) 不是单射。但 \(\rho\) 是满射：每个辛矩阵可由 Dehn 扭转实现。因此 \(\mathrm{Sp}(2g,\mathbb Z)\) 是 \(\mathrm{MCG}(\Sigma_g)\) 的同调商。

\emph{对 Ising 模型的物理后果。}曲面 \(\Sigma_g\) 上的 Ising 配分函数 \(Z_{\Sigma_g}\) 在微分同胚 \(\phi\in\mathrm{Diff}^+(\Sigma_g)\) 下不变（局部 Hamiltonian 只依赖邻域结构），因此映射类群作用平凡。但扭转配分函数 \(Z^{\boldsymbol\upsilon}\) 在 \(\phi\) 下变换为 \(Z^{\rho(\phi)\boldsymbol\upsilon}\)——扭转 character 被辛表示拉回。这给出约束：\(2^{2g}\) 个扭转配分函数在 \(\mathrm{Sp}(2g,\mathbb Z)\) 作用下置换。环面 \(g=1\) 时 \(\mathrm{Sp}(2,\mathbb Z)=\mathrm{SL}(2,\mathbb Z)\) 作用于四个扭转 \(\F_2^2\) 上，保持 Arf 不变量（偶/奇 spin structure 分类），与 Kac--Ward Pfaffian 的符号行为一致。

\emph{模空间与 Ising 能谱。}\(\mathrm{MCG}(\Sigma_g)\) 是 Teichm\"uller 空间 \(\mathcal T_g\) 的覆叠群，商 \(\mathcal T_g/\mathrm{MCG}(\Sigma_g)=\mathcal M_g\) 是模空间。Ising 能谱在 \(\mathcal T_g\) 上定义，在 \(\mathrm{MCG}\) 作用下不变，因此降为 \(\mathcal M_g\) 上的函数。辛表示 \(\rho\) 描述能谱在模空间上的单值性——沿模空间非平凡回路一周，扭转扇区按 \(\rho\) 混合。这是拓扑量子场论中 modular functor 的经典原型。辛表示的核与像满足
\begin{equation*}
\ker\rho=\mathrm{I}(\Sigma_g)\ (\text{Torelli 群}),\qquad \operatorname{im}\rho=\mathrm{Sp}(2g,\mathbb Z).
\end{equation*}







\end{derivation}

\emph{[曲面分类定理与 Van Kampen 基本群]}
亏格 \(g\) 完全分类闭可定向曲面：\(\Sigma_g\) 与 \(\Sigma_{g'}\) 同胚当且仅当 \(g=g'\)。基本群 \(\pi_1(\Sigma_g)\) 有标准展示，其 Abel 化给出 \(\mathrm H_1\)。分三步建立。

\emph{第一步：分类定理的证明骨架。}闭可定向曲面 \(\Sigma\) 有三角剖分（Rado 定理）。对剖分做如下化简：(i) 若两个三角形共享一条边且并集是圆盘，合并为一个面，减少面数；(ii) 若一条边只属于一个面（"悬挂边"），删去它和该面，减少边数和面数；(iii) 重复直到每个面边界是 \(4g\) 边形。最终得到标准多边形表示：单个 \(4g\) 边形，边界按
\begin{equation*}
  a_1b_1a_1^{-1}b_1^{-1}a_2b_2a_2^{-1}b_2^{-1}\cdots a_gb_ga_g^{-1}b_g^{-1}
\end{equation*}
粘合。亏格 \(g\) 由 Euler 特征 \(\chi=2-2g\) 唯一确定，因此同胚的曲面有相同 \(g\)。反之，相同 \(g\) 给出相同标准多边形，粘合后同胚。

\emph{第二步：Van Kampen 给出基本群。}标准 \(4g\) 边形 \(P_g\) 是可缩的（圆盘），\(\pi_1(P_g)=\{1\}\)。粘合边界边后，\(\Sigma_g\) 是 \(P_g\) 的商空间。由 Van Kampen 定理（或等价地，CW 复形的基本群公式），\(\pi_1(\Sigma_g)\) 由边界边的代表 loop \(a_1,b_1,\ldots,a_g,b_g\) 生成，唯一关系来自边界遍历：
\begin{equation*}
  \pi_1(\Sigma_g)=\langle a_1,b_1,\ldots,a_g,b_g\mid[a_1,b_1][a_2,b_2]\cdots[a_g,b_g]=1\rangle,
\end{equation*}
其中 \([a_i,b_i]=a_ib_ia_i^{-1}b_i^{-1}\) 是换位子。球面 \(g=0\) 给出 \(\pi_1(S^2)=\{1\}\)（空生成元集，空关系）。环面 \(g=1\) 给出 \(\pi_1(T^2)=\mathbb Z^2=\langle a,b\mid ab=ba\rangle\)（单一换位子关系即交换性）。

\emph{第三步：Abel 化与同调。}\(\mathrm H_1(\Sigma_g;\mathbb Z)\) 是 \(\pi_1(\Sigma_g)\) 的 Abel 化。换位子 \([a_i,b_i]\) 在 Abel 化中恒为 \(1\)，因此唯一关系消失，得到
\begin{equation*}
  \mathrm H_1(\Sigma_g;\mathbb Z)\cong\mathbb Z^{2g}=\mathbb Z a_1\oplus\mathbb Z b_1\oplus\cdots\oplus\mathbb Z a_g\oplus\mathbb Z b_g.
\end{equation*}
模二化给出 \(\mathrm H_1(\Sigma_g;\F_2)\cong\F_2^{2g}\)，与\eqnrefs{eq:phys-sm-genus-h1-dimension}一致。高阶同调由 Poincaré 对偶给出：\(\mathrm H_0\cong\mathbb Z\)（连通），\(\mathrm H_2\cong\mathbb Z\)（可定向闭曲面），\(\mathrm H_k=0\)（\(k>2\)）。

\emph{覆叠空间与 Ising 的周期边界。}\(\Sigma_g\) 的万有覆叠是双曲平面 \(\mathbb H^2\)（\(g\ge2\)）、欧氏平面 \(\mathbb R^2\)（\(g=1\)）或球面 \(S^2\)（\(g=0\)），覆叠群为 \(\pi_1(\Sigma_g)\)。Ising 模型在 \(\Sigma_g\) 上的周期边界条件等价于在万有覆叠上定义、再按 \(\pi_1\) 商化。\(g\ge2\) 时覆叠群是非 Abel 自由群（模单一关系），使扭转 character 的分类比环面（Abel \(\mathbb Z^2\)）复杂——但模二化后仍归为 \(\F_2^{2g}\) 的线性代数，这是同调比基本群更适合 Ising 模型的根本原因。

\emph{物理应用。}在拓扑量子计算中，亏格 \(g\) 曲面上的 Ising 任意子基态空间维数为 \(2^{g-1}(2+1)^g\)（Ising TQFT），与 \(\F_2^{2g}\) 的 \(4^g\) 不同——前者含任意子类型的额外简并。这一差异来自 TQFT 的非 Abel 性，是拓扑量子纠错码设计的数学基础。

\emph{分类定理的推广。}不可定向曲面由亏格 \(g\) 和不可定向性标记分类（如 Klein 瓶 \(g=1\)），其基本群含扭转元素。Ising 模型在不可定向曲面上的配分函数需要额外考虑 Pin 结构,结构（自旋结构的推广），与可定向情形有本质差异——例如 Klein 瓶上的 Ising 模型无自发磁化，因为不可定向性破坏了自旋的全球一致取向。

\emph{同调 vs 基本群的物理选择。}Ising 模型的周期边界条件用同调 \(\mathrm H_1(\Sigma_g;\F_2)\) 而非基本群 \(\pi_1(\Sigma_g)\) 分类，根本原因是 Ising 自旋 \(\mathbb Z_2\) 对称性使扭转 character 只取 \(\pm1\)，模二化后非 Abel 的 \(\pi_1\) 退化为 Abel 的 \(\mathrm H_1\)。对连续对称性模型（如 XY 模型 \(\mathrm U(1)\)），需用 \(\pi_1\) 的完整 \(U(1)\) 表示分类，同调不再充分。

\emph{计算复杂度与算法意义。}曲面亏格的计算可在 \(O(n)\) 时间完成（\(n\) 是三角剖分面数），通过 Euler 特征 \(\chi=V-E+F=2-2g\) 直接给出。但基本群 \(\pi_1(\Sigma_g)\) 的字问题（判断一个字是否代表单位元）在 \(g\ge2\) 时是非平凡的——虽然可解（Dehn 算法，\(O(n)\)），但群的复杂性使许多关于曲面上 Ising 模型的计算问题（如精确配分函数在亏格 \(g\) 曲面上的闭式）在 \(g\) 增大时变得困难。这一计算复杂度阶梯反映了拓扑复杂度与物理可解性之间的张力。







\emph{[Gauss--Bonnet 定理与 Euler 特征的曲率积分]}
Gauss--Bonnet 定理把曲面的局部几何（曲率）与全局拓扑（Euler 特征）联系起来，是微分几何与拓扑学的核心桥梁。分四步建立。

\emph{第一步：局部 Gauss--Bonnet。}对可定向曲面 \(\Sigma\) 上的测地三角形 \(T\)（三边为测地线），Gauss 曲率 \(K\) 在 \(T\) 上的积分与三角形内角 \(\alpha_1,\alpha_2,\alpha_3\) 满足
\begin{equation*}
  \int_T K\,\dd A=\alpha_1+\alpha_2+\alpha_3-\pi.
\end{equation*}
证明：对测地三角形 \(T\)，Stokes 定理给出 \(\int_T K\,\dd A+\int_{\partial T}k_g\,\dd s+\sum(\pi-\alpha_i)=2\pi\chi(T)=2\pi\)。测地线 \(k_g=0\)，\(\chi(T)=1\)（圆盘），故 \(\int_T K\,\dd A=\sum\alpha_i-\pi\)。

\emph{第二步：全局 Gauss--Bonnet。}对闭可定向曲面 \(\Sigma_g\)，三角剖分为 \(\{T_j\}\)，求和
\begin{equation*}
  \int_{\Sigma_g}K\,\dd A=\sum_j\int_{T_j}K\,\dd A=\sum_j(\alpha_{j1}+\alpha_{j2}+\alpha_{j3}-\pi).
\end{equation*}
每个顶点处三角形内角和为 \(2\pi\)（共 \(V\) 个顶点），每面贡献 \(-\pi\)（共 \(F\) 个面），每条边被两个三角形共享：
\begin{equation*}
  \int_{\Sigma_g}K\,\dd A=2\pi V-\pi F=2\pi(V-\frac{F}{2}).
\end{equation*}
由 Euler 公式 \(V-E+F=2-2g\) 和 \(3F=2E\)（三角剖分），
\begin{equation}
  \int_{\Sigma_g}K\,\dd A=2\pi(V-E+F)=2\pi\chi(\Sigma_g)=4\pi(1-g).
  \label{eq:phys-sm-gauss-bonnet}
\end{equation}

\emph{第三步：带边界推广。}对带边界曲面 \(\Sigma\)（边界 \(\partial\Sigma\) 由 \(n\) 条分段光滑闭曲线组成，外角 \(\theta_j\)），
\begin{equation*}
  \int_\Sigma K\,\dd A+\int_{\partial\Sigma}k_g\,\dd s+\sum_j\theta_j=2\pi\chi(\Sigma).
\end{equation*}
对圆盘（\(\chi=1\)，\(n=1\)，\(\theta=0\)）：\(\int K\,\dd A+\oint k_g\,\dd s=2\pi\)，给出平面曲线全曲率 \(\oint k_g\,\dd s=2\pi\)（Hopf Umlaufsatz）。对环面去掉一个圆盘（\(\chi=0\)）：\(\int K\,\dd A+\oint k_g\,\dd s=0\)。

\emph{第四步：物理应用。}(i) Ising 模型在曲面上的临界行为：Gauss--Bonnet 把亏格 \(g\)（通过 \(\chi\))与曲率积分联系，连续曲面 Ising 的有限尺寸修正确含 \(\int K\,\dd A=4\pi(1-g)\) 项。(ii) 拓扑场论：Chern--Gauss--Bonnet \(\int\operatorname{tr}(F\wedge F)=4\pi^2 c_2(E)\) 把规范场陈数与曲率积分联系，是拓扑绝缘体和量子 Hall 效应的数学基础。(iii) 曲面上随机游走：返回概率与 \(\chi\) 相关，Gauss--Bonnet 把"曲率正则返回快"（球面 \(g=0\)）与"曲率负则返回慢"（\(g\ge2\)）的直觉精确化。(iv) 共形几何：在共形变换 \(g\to\ee^{2\phi}g\) 下，\(\int K\,\dd A\) 变为 \(\int(K-\Delta\phi)\,\dd A\)，Laplacian 项积分为零（闭曲面），故 Euler 特征是共形不变量——这是 Uniformization 定理（闭曲面由亏格和共形类分类）的积分基础。






